\documentclass[11pt,a4paper]{article}
\usepackage[T1]{fontenc}
\usepackage[utf8]{inputenc}
\usepackage[margin=1in]{geometry}
\usepackage{amssymb,amsmath}
\usepackage[protrusion=true,expansion=false]{microtype}
\usepackage{enumitem}
\setlist{nosep,leftmargin=1.5em}
\usepackage{parskip}
\usepackage{hyperref}
\hypersetup{colorlinks=true,linkcolor=black,urlcolor=blue}
\newcommand{\euro}{EUR}
\pagestyle{plain}
\begin{document}
\section*{Word-count declaration}

Manuscript: \emph{A Comparability Protocol for Benchmarking Relational Database Access Layers in Express.js} (submission build \texttt{ist\_main.tex}).

The journal's Guide for Authors states a \textbf{15,000-word maximum for research papers}, that \textbf{references and appendixes count against the total}, and that \textbf{figures and tables count 200 words each}. This declaration uses a deliberately conservative implementation of that rule.

\begin{center}\begin{tabular}{ll}
\hline
Component & Count \\
\hline
Seven body sections (\texttt{texcount} sum, including section headings and caption text encountered in those files) & 9,721 \\
Structured abstract (\texttt{texcount} sum) & 295 \\
Remaining front matter and required declarations (title/keywords, CRediT, competing interest, data availability, funding, and AI-use disclosures) & 442 \\
Main-text tables and figures (2 x 200) & 400 \\
Rendered reference list (51 entries; \texttt{pdftotext} count from the first entry to the end, form feeds stripped) & 1,536 \\
\textbf{Conservative IST submission equivalent} & \textbf{12,394} \\
\hline
\end{tabular}\end{center}

The controlling total is therefore \textbf{12,394 words}, leaving \textbf{2,606 words of headroom} under the 15,000-word limit. The structured abstract contains 295 words under the \texttt{texcount} sum convention, below the 300-word limit. The manuscript has two floats, Figure 1 and Table 1, so the float allowance is 400 words.

The decomposition is reproducible in four commands, so a reader need not take the total on trust: \texttt{texcount -1 -sum -q paper/sections/*.tex} gives the body; \texttt{texcount -1 -sum -q paper/ist/ist\_main.tex} gives the front matter, abstract and declarations together, from which the abstract is subtracted to avoid counting it twice; the float allowance is 2 x 200; and the reference list is counted from the built PDF with \texttt{pdftotext ... | tr -d '\textbackslash\{\}f' | awk '/\textasciicircum{}References\$/\{f=1;next\} f' | wc -w}.

\section*{Revision against the previous declaration}

The previous declaration reported 14,992 words. This revision removes \textbf{2,598 words}, a 17\textbackslash\{\}\% reduction, in response to the editorial office's length notice. The manuscript is now 43 pages, down from 52.

Reductions by component:

\begin{center}\begin{tabular}{llll}
\hline
Component & Previous & Now & Change \\
\hline
Introduction & 1,085 & 899 & -186 \\
Related work & 1,366 & 1,157 & -209 \\
Study Design & 4,303 & 3,311 & -992 \\
Results & 2,072 & 1,700 & -372 \\
Discussion & 1,080 & 888 & -192 \\
Threats to Validity & 1,606 & 1,361 & -245 \\
Conclusion & 353 & 345 & -8 \\
References & 1,991 & 1,536 & -455 \\
\hline
\end{tabular}\end{center}

Three mechanisms were used, in this order of preference. First, supporting detail was \textbf{relocated to the online supplement}, which gained eight sections carrying it; no content was lost, only moved. Second, material \textbf{duplicated across sections} was stated once, chiefly where the Introduction restated the protocol specification given normatively in Study Design and the prior-work enumeration given in Related Work. Third, \textbf{fourteen references} peripheral to the argument were removed, nine supporting context rather than a claim and five duplicating a single point about ORM query anti-patterns; every claim they supported remains, and no cited entry is missing from the bibliography.

No measured value, research question, access layer, access pattern, or engine was removed. Every limitation and threat to validity remains stated in the article; only the measurements supporting them moved to the supplement.

\section*{Note on the counting method}

The editorial office reported 17,581 words for the previous submission, against the 14,992 declared above. The authors could not reproduce that figure by any counting method applied to the submitted files: the complete extracted text of the previous PDF, including the reference list, was 15,134 words, and the most conservative variant, adding the 200-word float allowance to that total, was 15,534. The present revision reduces the extracted-text total to 12,547 words, so it falls below 15,000 under both a proportional and an additive model of the difference. Should the total still exceed the limit under the office's method, the authors would be grateful for the counting convention used, and will reduce further accordingly.

Count date: 2 August 2026.

\end{document}
